\newpage
\pdfbookmark[1]{\infoname}{info}

% Eesti keeles info.
\begin{info}
\begin{abstract}
Käesolev bakalaureusetöö uurib tehisintellekti (TI) rakendamise teostatavust eestikeelse tehnilise informaatika lõputöö genereerimisel. Töös võrreldakse nelja juhtivat suurt keelemudelit koos vastavate agentsete tööriistadega -- Google Gemini 3.1 Pro (Antigravity), OpenAI GPT-5.4 (ChatGPT Codex laiendus Visual Studio Code'is), Anthropic Claude Sonnet 4.6 (Cowork) ning Anthropic Claude Opus 4.6 (Cowork) -- nende võimekuse osas luua struktureeritud, akadeemiliselt korrektset ja sisuliselt väärtuslikku eestikeelset teksti. Kõigi kombinatsioonide puhul kasutati ühte ja sama lühikest eestikeelset viibet, mis palus luua LaTeX-vormingus terviklik bakalaureusetöö, mis vastab Tartu Ülikooli arvutiteaduse instituudi nõuetele, valida ise informaatikaga seotud teema ning kasutada töökausta paigutatud UniTartuCS malli. Cowork on sama ettevõtte toodetud Claude Code'ist eraldiseisev agentne tööriist. Genereeritud tekstide kvaliteeti hinnati struktuuri, keelelise korrektsuse (kirjavigade tiheduse), viidete hulga ja kontrollitavuse, eestikeelse akadeemilise stiili (anglitsismide ja esimese isiku asesõnade kasutuse) ning mahu põhjal. Eksperimendi tulemused paljastasid kombinatsioonide vahel olulisi erinevusi. Claude Opus 4.6 + Cowork genereeris 61-leheküljelise eestikeelse lõputöö paroolihalduri teemal. Claude Sonnet 4.6 + Cowork genereeris 40-leheküljelise eestikeelse lõputöö TI-põhise lõputöö genereerimise teemal. GPT-5.4 keeldus tavapärase ChatGPT vestlusliidese kaudu täielikku tööd kirjutamast, kuid Codex laiendusega õnnestus tal genereerida 47-leheküljeline eestikeelne töö RAG-põhise küsimus-vastus süsteemi teemal. Gemini 3.1 Pro + Antigravity genereeris 39-leheküljelise eestikeelse töö generatiivse TI mõjust programmeerimisele, mille empiiriline osa oli täielikult fabritseeritud. Tulemused näitavad, et mudel-tööriista kombinatsioonide suutlikkus erineb drastiliselt nii ülesandest arusaamise, keelevaliku, ortograafilise täpsuse kui ka juhiste järgimise osas. Töö annab soovitused TI tööriistade efektiivseks kasutamiseks lõputöö kirjutamise protsessis ning kaardistab valdkonnad, kus inimese panus on endiselt asendamatu.
\end{abstract}

\keywords{tehisintellekt, keelemudelid, lõputöö genereerimine, eesti keel, teostatavusuuring, Cowork, Claude Code}

\cercs{P170 Arvutiteadus, arvutusmeetodid, süsteemid, juhtimine (automaatjuhtimisteooria)}
\end{info}


% Inglise keeles info.
\begin{otherInfo}{english}{Feasibility Study: Generating an Estonian-Language Technical Computer Science Thesis Using Artificial Intelligence}
\begin{abstract}
This bachelor's thesis investigates the feasibility of applying artificial intelligence (AI) to generate an Estonian-language technical computer science thesis. The study compares four leading large language models paired with corresponding agentic tools -- Google Gemini 3.1 Pro (Antigravity), OpenAI GPT-5.4 (ChatGPT chat interface and ChatGPT Codex VS Code extension), Anthropic Claude Sonnet 4.6 (Cowork), and Anthropic Claude Opus 4.6 (Cowork) -- in terms of their ability to produce structured, academically correct, and substantively valuable Estonian-language text. A single short Estonian-language user prompt was used for every combination, instructing the model to produce a complete LaTeX-formatted bachelor's thesis meeting the requirements of the University of Tartu's Institute of Computer Science, to choose a computer-science-related topic itself, and to use the UniTartuCS LaTeX template placed in the working directory. Cowork is a separate Anthropic product from Claude Code. Output quality was assessed using structure, orthographic accuracy, the number and verifiability of citations, Estonian academic style (anglicism and first-person pronoun usage), and document length. The experiment revealed significant differences between combinations. Claude Opus 4.6 + Cowork produced a 61-page Estonian-language thesis on a password manager. Claude Sonnet 4.6 + Cowork produced a 40-page Estonian-language thesis on AI-based thesis generation. GPT-5.4 declined to write a full thesis via the standard ChatGPT chat interface, but via the Codex extension succeeded in generating a 47-page Estonian-language thesis on a RAG-based question-answering system. Gemini 3.1 Pro + Antigravity generated a 39-page Estonian-language thesis on the impact of generative AI on programming, whose empirical part was wholly fabricated. Results demonstrate that model-tool combinations differ drastically in task comprehension, language selection, orthographic accuracy, and instruction adherence. The thesis provides recommendations for the effective use of AI tools in the thesis writing process and maps areas where human contribution remains irreplaceable.
\end{abstract}

\keywords{artificial intelligence, language models, thesis generation, Estonian language, feasibility study, Cowork, Claude Code}

\cercs{P170 Computer science, numerical analysis, systems, control}
\end{otherInfo}
